\documentclass[11pt]{article}

\usepackage[margin=2.8cm]{geometry}
\usepackage{amsmath,amssymb,amsthm}
\usepackage{stmaryrd}
\usepackage{enumitem}
\usepackage[hidelinks]{hyperref}

\title{A Formalization of the Finite-Trace Calculus of
  ``Automated Verification of Finite-trace Properties for Recursive
  Programs''\\[0.3em]
  \large How the formalization differs from the thesis}
\author{}
\date{}

% ---------------------------------------------------------------------------
% Notation

\newcommand{\sem}[1]{[\![ #1 ]\!]}
\newcommand{\fsem}[1]{|\!|#1|\!|}
\newcommand{\Str}{{\mathcal{S}_{\mathit{tr}}}}
\newcommand{\Seff}{{\mathcal{S}_{\mathit{eff}}}}
\newcommand{\chop}{\smallfrown}
\newcommand{\Id}{\mathsf{Id}}
\newcommand{\Sb}[2]{\mathsf{Sb}^{#1}_{#2}}
\newcommand{\dup}{\sharp}
\newcommand{\State}{\mathit{State}}
\newcommand{\kw}[1]{\mathbf{#1}}
\newcommand{\skipS}{\kw{skip}}
\newcommand{\ite}[3]{\kw{if}\ #1\ \kw{then}\ #2\ \kw{else}\ #3}
\newcommand{\call}[1]{#1()}
\newcommand{\Tr}{\State^{*}}
\newcommand{\Pre}{\mathit{Pre}}
\newcommand{\Post}{\mathit{Post}}
\newcommand{\tc}{\mathrm{tc}}
\newcommand{\ftc}{\mathrm{ftc}}
\newcommand{\pred}{\mathrm{pred}}
\newcommand{\distr}{\mathrm{distr}}
\newcommand{\norm}{\mathrm{norm}}
\newcommand{\seqn}{\mathrm{seqn}}
\newcommand{\first}{\mathrm{first}}
\newcommand{\last}{\mathrm{last}}
\newcommand{\tl}{\mathrm{tail}}
\newcommand{\B}{\mathcal{B}}
\newcommand{\A}{\mathcal{A}}
\newcommand{\eff}{\rhd}
% A numbered item of the thesis [FS], e.g. \Th{Definition~4.8}.
\newcommand{\Th}[1]{#1 of FS}

\theoremstyle{plain}
\newtheorem{theorem}{Theorem}[section]
\newtheorem{lemma}[theorem]{Lemma}
\newtheorem{proposition}[theorem]{Proposition}
\newtheorem{corollary}[theorem]{Corollary}
\theoremstyle{definition}
\newtheorem{definition}[theorem]{Definition}
\theoremstyle{remark}
\newtheorem{remark}[theorem]{Remark}

\begin{document}
\maketitle

\section{Introduction}

This note accompanies a HOL4 formalization of the soundness results of
Chapter~4 of the master's thesis \emph{Automated Verification of
Finite-trace Properties for Recursive Programs} by
Stevanovic~\cite{FS} (henceforth FS). The thesis reduces the verification
of a Rec program against a trace formula to a validity problem: the trace
constraints $\tc(S)$ describe the traces of $S$, with procedure calls
abstracted by contracts, and the finite-trace calculus $\ftc(S, \phi)$
encodes the specification $\phi$. The formalization covers
\begin{itemize}
\item \Th{Lemma~4.16}: every effective trace satisfies the trace
  constraints;
\item \Th{Theorem~4.29}: a trace whose effective trace satisfies
  $\ftc(S, \phi)$ is in $\fsem{\phi}$;
\item \Th{Lemma~4.30}: validity of $\tc(S) \Rightarrow \ftc(S, \phi)$
  implies $\Str\sem{S} \subseteq \fsem{\phi}$;
\end{itemize}
and examples of the thesis that check that the definitions and results
apply as intended (Section~\ref{sec:examples}), and that show that the
converse of \Th{Lemma~4.30} fails (Section~\ref{sec:incomplete}).
It is built on the formalization of the paper by Gurov and
H\"ahnle~\cite{GH}, whose language, semantics and trace logic the thesis
uses (\Th{Chapter~2}). The note on that formalization describes its
conventions; this note relies on them.

The results hold, but several definitions of the thesis do not support its
proofs as stated; with the definitions as stated, \Th{Lemma~4.16},
\Th{Theorem~4.29} and \Th{Lemma~4.30} fail. The noteworthy changes are:
\begin{itemize}
\item \emph{Effective traces.} \Th{Definition~4.8} abstracts only a
  procedure call at the top level of a statement, but the proofs treat every
  call of the statement this way. With the definition as stated,
  \Th{Lemma~4.16} fails (Proposition~\ref{prop:seff}). The formalization
  defines effective traces compositionally, as a relation between a trace
  and its effective trace (Definition~\ref{def:eff}). This also makes the
  bijection of \Th{Lemma~4.9}, and the determinism of Rec that it relies
  on, unnecessary.
\item \emph{Normal form of statements.} The shape condition that
  \Th{Lemma~4.16} assumes (no conditional as the left operand of a
  sequence) is too weak, because $\tc$ re-associates sequences and can
  create such a conditional (Proposition~\ref{prop:shape}). The
  formalization requires a stronger normal form
  (Definition~\ref{def:nf}) and normalizes statements and procedure bodies
  with a corrected $\norm$ (Definition~\ref{def:norm}). As a result,
  \Th{Lemma~4.30} holds for every statement without assumptions about its
  shape (Theorem~\ref{thm:vc}).
\item \emph{Chop and conjunction.} The second equation of \Th{Lemma~2.22},
  $(A \cap B)\chop C = (A\chop C) \cap (B\chop C)$, is false; only the
  inclusion from left to right holds. Hence $\distr$ (\Th{Definition~4.23})
  does not preserve the semantics (\Th{Lemma~4.24} fails,
  Proposition~\ref{prop:distr}). The proof of \Th{Theorem~4.29} uses
  \Th{Lemma~4.24} for formulas $(\mu X.\psi)\chop\phi_2$, and with the
  thesis's $\distr$, \Th{Theorem~4.29} and \Th{Lemma~4.30} fail
  (Theorem~\ref{thm:ftc-thesis}). The formalization does not distribute
  chop over conjunction (Definition~\ref{def:distr}).
\item \emph{Distribution and grouping of chops.} The thesis writes chop
  sequences without brackets, and \Th{Example~4.25} regroups them
  implicitly. \Th{Examples~4.31 and~4.32} moreover need chop distributed
  over disjunction on the right, which \Th{Definition~4.23} does not do:
  as stated, the case for a call in \Th{Definition~4.28} does not apply
  to them. The corrected $\distr$ does both, and the formalization verifies
  these examples (Theorem~\ref{thm:ex431}).
\item \emph{Fixed points of calls.} \Th{Definition~4.28} writes the
  specification of a call as $\mu X_m.\psi$, which suggests that the
  recursion variable must belong to the called procedure. The thesis's
  \Th{Examples~5.8 to~5.12} instead let a call consume any fixed point,
  such as the one that $\mathrm{expand}$ introduces for a state formula.
  The formalization accepts any fixed point (Definition~\ref{def:ftc}).
\item \emph{Partial functions.} $\pred$ and $\ftc$
  (\Th{Definitions~4.17 and~4.28}) need not terminate. The formalization
  gives them fuel and makes them false when it runs out. Theorem~4.29 then
  needs neither the hypothesis that $\ftc(S, \phi)$ is well-defined nor that
  $\phi$ is normalised (Theorem~\ref{thm:ftc}). Nor does it need the
  function $\mathrm{expand}$ (\Th{Definition~4.20}); it is formalized with
  \Th{Lemma~4.21} and normalises the specifications of the examples.
\item \emph{Proof of \Th{Theorem~4.29}.} The case of a recursive call
  relies on a side condition generated earlier in the recursion. The
  formalization makes this precise with an invariant on the set $X$ of
  side conditions, and proves the case of a first call by an inner
  induction on the length of the traces of the procedure body
  (Lemma~\ref{lem:side}).
\end{itemize}
The encodings are embedded shallowly here, as predicates on state
sequences (Section~\ref{sec:conv}). The first-order logic of
\Th{Section~4.1} is formalized on top of them, as described in a separate
note (\texttt{doc/fol.tex}).

Each numbered definition and result in this note corresponds to a HOL4
definition or theorem in the directory \texttt{hol/verif}. The file
\texttt{README.md} lists them by the numbers used here. References such as
``\Th{Definition~4.8}'' are to the thesis.

\section{Conventions}\label{sec:conv}

The language, the trace semantics $\Str$ and the semantics $\fsem{\phi}$ of
trace formulas are those of the formalization of GH. In particular, states
are total functions from variable names to integers, procedure tables are
lists looked up by first match, and an undeclared procedure has no traces.
Statement variables are not part of the thesis's language; they are
interpreted by $\emptyset$, so they have no traces. Substitution
$\phi[\psi/X]$ is plain substitution; all formulas below are closed when
they are substituted, so this agrees with capture-avoiding substitution.
As in GH's formalization, the recursion variable $X_m$ of a procedure $m$
is the name $m$ itself.

\paragraph{State sequences.} The thesis's formulas over state variables
$s^{(1)}, s^{(2)}, \dots$ are embedded shallowly: a \emph{state sequence} is
a function $g : \mathbb{N} \to \State$, with $g(i)$ in place of $s^{(i)}$,
and an encoding is a predicate on state sequences. A trace
$u = u_0 \cdots u_{k}$ is \emph{placed at $i$ in $g$} if $g(i+j) = u_j$ for
all $j \leq k$; positions after the trace are arbitrary, which is the
thesis's padding (\Th{Example~4.3}). The relations $\Id$ and $\Sb{a}{x}$ of
\Th{Definition~4.1} become equalities of states,
$g(i+1) = g(i)$ and $g(i+1) = g(i)[x \mapsto \A\sem{a}(g(i))]$, and
$\B\sem{b}(s^{(i)})$ becomes $\B\sem{b}(g(i))$. Since states are total,
this is the thesis's structure $M$ with all variables in $xs$. A link to a
first-order logic with finitely many variables, as needed for SMT solving,
is left for later.

\section{Contracts and effective traces}

\begin{definition}[\Th{Definitions~4.5 and~4.6}]
  A contract assigns to each procedure $m$ a precondition
  $\Pre_m \subseteq \State$ and a postcondition
  $\Post_m \subseteq \State \times \State$. Every procedure
  \emph{conforms} with its contract if $\Pre_m(\first\,\sigma)$ implies
  $\Post_m(\first\,\sigma, \last\,\sigma)$ for every $m$ and every
  $\sigma \in \Str\sem{\call{m}}$.
\end{definition}

\Th{Definition~4.8} sets
$\Seff\sem{\call{m}} = \{s_1 s_2 s_n \mid s_1 s_2 \cdots s_n \in
\Str\sem{\call{m}}\}$ and $\Seff\sem{S} = \Str\sem{S}$ for every other
statement. A call inside a sequence or a conditional is therefore not
abstracted. The proofs of \Th{Lemma~4.16 and Theorem~4.29} use effective
traces in which every call of the statement is abstracted, as in the
proof of \Th{Lemma~4.16} for $\call{m}; S_2$. With the definition as
stated, the lemma fails.

\begin{proposition}\label{prop:seff}
  Let $T = \{m\{x := 1; x := 2\}\}$ with the contract
  $\langle \mathit{true},\ x = 2\rangle$, which $m$ conforms with. Let
  $S = \skipS; \call{m}$, let $s$ be the state that maps every variable
  to $0$, and let $s_1 = s[x \mapsto 1]$ and $s_2 = s_1[x \mapsto 2]$. Then
  $u = s\,s\,s\,s_1\,s_2 \in \Seff\sem{S}$ as in \Th{Definition~4.8},
  but if $u$ is placed at $1$ in $g$, then $g$ does not satisfy $\tc_1(S)$:
  the conjunct $\Post_m(g(3), g(4))$ requires $g(4)(x) = 2$, but
  $g(4) = s_1$.
\end{proposition}

The formalization relates each trace to its effective trace directly.

\begin{definition}[Effective traces, corrected \Th{Definition~4.8}]
  \label{def:eff}
  The relation $\sigma \eff_S \tau$ ($\tau$ is an effective trace of the
  trace $\sigma$ of $S$) is defined inductively:
  \begin{gather*}
    s\,s \eff_{\skipS} s\,s \qquad
    s\,s' \eff_{x := a} s\,s' \text{ where } s' = s[x \mapsto \A\sem{a}(s)]
    \qquad
    \frac{\sigma \in \Str\sem{\call{m}}}
         {\sigma \eff_{\call{m}} \first\,\sigma\ \first\,\sigma\ \last\,\sigma}
    \\[0.5em]
    \frac{\sigma_1 \eff_{S_1} \tau_1 \quad \sigma_2 \eff_{S_2} \tau_2 \quad
          \last\,\sigma_1 = \first\,\sigma_2}
         {\sigma_1 \cdot \tl\,\sigma_2 \eff_{S_1; S_2}
          \tau_1 \cdot \tl\,\tau_2}
    \\[0.5em]
    \frac{\B\sem{b}(s) \quad s\,\sigma \eff_{S_1} s\,\tau}
         {s\,s\,\sigma \eff_{\ite{b}{S_1}{S_2}} s\,s\,\tau}
    \qquad
    \frac{\lnot\B\sem{b}(s) \quad s\,\sigma \eff_{S_2} s\,\tau}
         {s\,s\,\sigma \eff_{\ite{b}{S_1}{S_2}} s\,s\,\tau}
  \end{gather*}
\end{definition}

The body of a called procedure is not abstracted further, as in the
thesis. If $\sigma \eff_S \tau$, then $\sigma \in \Str\sem{S}$, both traces
have at least two states, and they have the same first and the same last
state.

\begin{lemma}[replaces \Th{Lemma~4.9}]\label{lem:eff-exists}
  Every $\sigma \in \Str\sem{S}$ has an effective trace $\tau$ with
  $\sigma \eff_S \tau$.
\end{lemma}

The proof is by induction on $S$. The bijection of \Th{Lemma~4.9} is not
needed: Theorem~\ref{thm:ftc} is stated for a trace together with an
effective trace of it, and its proof constructs the effective traces of
subtraces from Lemma~\ref{lem:eff-exists} where they are needed. So the
determinism of Rec is not used either.

\section{Trace constraints}

\begin{definition}[Normal form]\label{def:nf}
  A statement is \emph{atomic} if it is $\skipS$, an assignment, a call or
  a statement variable. A statement is in \emph{normal form} if the left
  operand of each of its sequences is atomic.
\end{definition}

So in normal form, sequences are associated to the right and have no
conditional as left operand.

\begin{definition}[\Th{Definition~4.13}]
  $\tc_i(S)$ is defined by the equations of \Th{Definition~4.13}, with
  $\tc_i(Y) = \mathit{true}$ for a statement variable $Y$. In particular,
  \[\tc_i((S_1; S_2); S_3) = \tc_i(S_1; (S_2; S_3)),\]
  and
  $\tc_i(S_1; S_2) = \tc_i(S_1) \land \tc_{i+1}(S_2)$ if $S_1$ is neither a
  call nor a sequence.
\end{definition}

The re-association makes the recursion terminate under a weight in which
the left operand of a sequence counts twice. The thesis assumes,
without loss of generality, that no conditional is the left operand of a
sequence (text after \Th{Lemma~4.11}), and that its $\norm$
(\Th{Definition~4.33}) establishes this. But the case for $S_1 = S'_1;S'_2$
re-associates, which can move a conditional to the left of a sequence.
Then the index $i+1$ for $S_2$ is wrong.

\begin{proposition}\label{prop:shape}
  Let $S = (\skipS; \ite{\mathit{true}}{x := 1}{\skipS}); x := 2$. No
  conditional in $S$ is the left operand of a sequence, and the thesis's
  $\norm$ (\Th{Definition~4.33}) leaves $S$ unchanged. Let $s$ be the
  state that maps every variable to $0$, $s_1 = s[x \mapsto 1]$ and
  $s_2 = s_1[x \mapsto 2]$. Then $\sigma = s\,s\,s\,s_1\,s_2$ is a trace of
  $S$ and $\sigma \eff_S \sigma$, but if $\sigma$ is placed at $1$ in $g$,
  then $g$ does not satisfy $\tc_1(S)$: the conjunct for $x := 2$ is
  $g(4) = g(3)[x \mapsto 2]$, but $g(3) = s$ and $g(4) = s_1$.
\end{proposition}

With the stronger normal form, the lemma holds.

\begin{lemma}[\Th{Lemma~4.16}]\label{lem:tc}
  Assume that every procedure conforms with its contract. If $S$ is in
  normal form, $\sigma \eff_S \tau$, and $\tau$ is placed at $i$ in $g$,
  then $g$ satisfies $\tc_i(S)$.
\end{lemma}

The proof is by induction on $\sigma \eff_S \tau$. A statement variable
has no traces, so it does not occur as the left operand of a sequence in
the induction.

\begin{definition}[Normalization, corrected \Th{Definition~4.33}]
  \label{def:norm}
  Sequential composition in normal form is defined by
  \begin{align*}
    \seqn(S_1; S_2,\ S_3) &= S_1; \seqn(S_2, S_3)\\
    \seqn(\ite{b}{S_1}{S_2},\ S_3) &=
      \ite{b}{\seqn(S_1, S_3)}{\seqn(S_2, S_3)}\\
    \seqn(S_1, S_3) &= S_1; S_3 \quad \text{otherwise.}
  \end{align*}
  Then $\norm(S_1;S_2) = \seqn(\norm(S_1), \norm(S_2))$, $\norm$ commutes
  with conditionals, and it is the identity on atomic statements. A
  procedure table is normalized by normalizing every body.
\end{definition}

\begin{lemma}[generalizes \Th{Lemma~4.11}]\label{lem:norm}
  $\norm(S)$ is in normal form and has the same traces as $S$, also when
  the procedure table is normalized. More generally,
  $\Str'\sem{\seqn(S, S')}_{\rho,I} =
  \Str'\sem{S}_{\rho,I} \chop \Str'\sem{S'}_{\rho,I}$ for every
  interpretation $\rho$ of procedures.
\end{lemma}

The proof uses that chop is associative, distributes over union, and
commutes with restriction and duplication on the left. The
procedure-relative semantics $\Str'$ is that of GH's formalization; since
the transformer whose least fixed point interprets the procedures does
not change under normalization of the table, neither does $\Str$.

\section{The specification encoding}

\begin{definition}[\Th{Definition~4.17}]
  $\pred^n_i(\phi)$ is defined by the equations of \Th{Definition~4.17},
  with a fuel argument $n$: for a fixed-point formula,
  $\pred^0_i(\mu X.\phi) = \mathit{false}$ and
  $\pred^{n+1}_i(\mu X.\phi) = \pred^n_i(\phi[\mu X.\phi/X])$. The other
  cases keep $n$.
\end{definition}

The thesis notes that $\pred_i(\mu X.X)$ is not well-defined. With fuel,
$\pred$ is total, and the result below holds for any fuel.

\begin{lemma}[\Th{Lemma~4.19}]\label{lem:pred}
  If $\phi$ is closed and $g$ satisfies $\pred^n_i(\phi)$, then
  $g(i)\,g(i+1) \in \fsem{\phi}$.
\end{lemma}

The proof of \Th{Lemma~4.24} rests on \Th{Lemma~2.22}. Its second
equation fails: a trace in both $A\chop C$ and $B\chop C$ may be split at
different positions.

\begin{proposition}\label{prop:distr}
  $\fsem{(\Id \land \Id\chop\Id)\chop \mathit{true}} = \emptyset$, but
  \Th{Definition~4.23} distributes the chop to
  $(\Id\chop\mathit{true}) \land ((\Id\chop\Id)\chop\mathit{true})$, whose
  semantics contains $s\,s\,s$.
\end{proposition}

The proof of \Th{Theorem~4.29} uses \Th{Lemma~4.24} in the case for
$(\mu X.\psi)\chop\phi_2$, and that step is not valid for conjunctions.
Section~\ref{sec:distr-thesis} shows that \Th{Theorem~4.29} indeed fails
with the thesis's $\distr$.

\begin{definition}[\Th{Definition~4.20}]\label{def:expand}
  $\mathrm{expand}(\phi)$ is defined as in \Th{Definition~4.20}. Since
  $\mu X'.(\mathit{true} \lor \mathit{true}\chop X')$ is closed, the name
  $X'$ need not be fresh.
\end{definition}

\begin{lemma}[\Th{Lemma~4.21}]
  $\fsem{\mathrm{expand}(\phi)}_V = \fsem{\phi}_V$ for every $\phi$ and
  every valuation $V$. Moreover, $\mathrm{expand}(\phi)$ has the same free
  variables as $\phi$.
\end{lemma}

\begin{definition}[corrected \Th{Definition~4.23}]\label{def:distr}
  $\distr(\phi)$ is defined by
  \begin{align*}
    \distr((\psi_1 \lor \psi_2)\chop\phi_2) &=
      \distr(\psi_1\chop\phi_2) \lor \distr(\psi_2\chop\phi_2)\\
    \distr((\phi_1\chop\phi_2)\chop\phi_3) &=
      \distr(\phi_1\chop(\phi_2\chop\phi_3))\\
    \distr(\phi_1\chop\phi_2) &= \distr(\phi_1) \chop_d \distr(\phi_2)
      \quad\text{otherwise}
  \end{align*}
  and by applying $\distr$ to the subformulas of other formulas, where
  $\phi_1 \chop_d (\psi_1 \lor \psi_2) = (\phi_1\chop_d\psi_1) \lor
  (\phi_1\chop_d\psi_2)$ and $\phi_1 \chop_d \phi_2 = \phi_1\chop\phi_2$
  otherwise.
\end{definition}

Compared to the thesis, the case for conjunction is dropped. The second
equation makes explicit the regrouping of \Th{Example~4.25}, and
$\chop_d$ distributes chop over disjunction on the right, which
\Th{Examples~4.31 and~4.32} need (Section~\ref{sec:examples}). Without the
regrouping, $\distr$ leaves left-nested chops, for example from
$\mathrm{expand}$, and $\ftc$ is false on them, because $\pred$ is false on
a chop.

\begin{lemma}[\Th{Lemma~4.24}]
  $\fsem{\distr(\phi)}_V = \fsem{\phi}_V$ for every $\phi$ and every
  valuation $V$. Moreover, $\distr(\phi)$ has the same free variables as
  $\phi$.
\end{lemma}

The proof uses that chop is associative and distributes over union on
both sides. Together with \Th{Lemma~4.21}, this gives
\Th{Proposition~4.27}: $\fsem{\distr(\mathrm{expand}(\phi))} =
\fsem{\phi}$.

\section{The finite-trace calculus}

\begin{definition}[\Th{Definition~4.28}]\label{def:ftc}
  $\ftc^n_i(S, \phi, X)$ and the helper $\ftc'^n_i(S, \phi, X)$ are defined
  by the equations of \Th{Definition~4.28}, where $X$ is a list of pairs
  $(m, \mu Y.\psi)$ of a procedure and a fixed-point formula, and with fuel
  $n$: $\ftc^0_i(S, \phi, X) = \mathit{false}$,
  each case of $\ftc^{n+1}$ calls $\ftc^n$ or $\ftc'^n$, and each case of
  $\ftc'^n$ calls $\ftc^n$ and $\pred^n$. The case for a call is
  \[
    \ftc'^n_i(\call{m}, \phi_1\chop\mu Y.\psi, X) =
    \Pre_m(g(i+1)) \land \pred^n_i(\phi_1) \land C
  \]
  where $C$ is $\mathit{true}$ if $(m, \mu Y.\psi) \in X$, and otherwise the
  side condition
  \[
    \forall h.\ \tc_1(S_m)(h) \land \Pre_m(h(1)) \Rightarrow
      \ftc^n_1(S_m,\ \psi[\mu Y.\psi/Y],\ (m, \mu Y.\psi) :: X)(h)
  \]
  if $m\{S_m\} \in T$, and $\mathit{false}$ if $m$ is undeclared. All cases
  that the thesis leaves undefined, including statement variables, are
  $\mathit{false}$.
\end{definition}

\Th{Definition~4.28} writes the fixed point of a call as $\mu X_m.\psi$,
but a call may be specified by any fixed point $\mu Y.\psi$. The thesis's
own \Th{Examples~5.8 to~5.12} rely on this: a call consumes the fixed point
that $\mathrm{expand}$ introduces for a state formula (see the discussion
of \Th{Example~5.11}). Requiring $Y = X_m$ is sound but rejects these
examples. $X$ therefore records the whole fixed-point formula. The thesis
orders the cases by a lexicographic order
whose details it omits, and its \Th{Theorem~4.29} assumes that
$\ftc(S, \phi)$ is well-defined; with fuel, every $\ftc^n$ is a total
function and sound.

\begin{theorem}[\Th{Theorem~4.29}]\label{thm:ftc}
  Assume that every procedure conforms with its contract and every
  procedure body is in normal form. Let $\phi$ be closed, let
  $\sigma \eff_S \tau$, and let $\tau$ be placed at $i$ in $g$. If $g$
  satisfies $\ftc^n_i(S, \phi, [\,])$, then $\sigma \in \fsem{\phi}$.
\end{theorem}

The hypothesis that $\phi$ is normalised is dropped: the result holds for
any closed $\phi$. Normalization of $\phi$ only makes more formulas
provable.

The proof is by induction on $n$. The claim is generalized to arbitrary
lists $X$ that satisfy an invariant: $X$ is \emph{good below $L$} if for
every $(m, \mu Y.\psi) \in X$ with $m\{S_m\} \in T$, every
$\sigma' \in \Str\sem{S_m}$ with $|\sigma'| < L$ and
$\Pre_m(\first\,\sigma')$ is in $\fsem{\mu Y.\psi}$. The claim for fuel
$n$ is: if $\ftc^n_i(S, \phi, X)$ holds as in the theorem and $X$ is good
below $|\sigma|$, then $\sigma \in \fsem{\phi}$. The case of a recursive
call, $(m, \mu Y.\psi) \in X$, follows directly from the invariant, since the
trace of the body is shorter than the trace of the call. This is the
precise form of the thesis's appeal to the side condition ``in the
recursion that leads to this case''. The cases for sequences use that
the invariant is preserved for shorter traces, and that chop is
associative (for $\call{m}; S'$).

\begin{lemma}\label{lem:side}
  Assume the claim for fuel $n$, $m\{S_m\} \in T$, $\mu Y.\psi$ closed,
  $X$ good below $L$, and the side condition of
  Definition~\ref{def:ftc} for $m$, $\mu Y.\psi$, $X$ and fuel $n$. Then
  every $\sigma' \in \Str\sem{S_m}$ with $|\sigma'| < L$ and
  $\Pre_m(\first\,\sigma')$ is in $\fsem{\mu Y.\psi}$.
\end{lemma}

This is the case of a first call. The proof is by induction on the length
of $\sigma'$. Take an effective trace $\tau'$ of $\sigma'$
(Lemma~\ref{lem:eff-exists}) and place it at $1$ in some $h$. By
Lemma~\ref{lem:tc}, which needs the body $S_m$ in normal form, $h$
satisfies $\tc_1(S_m)$, so the side condition gives
$\ftc^n_1(S_m, \psi[\mu Y.\psi/Y], (m, \mu Y.\psi) :: X)$. The list
$(m, \mu Y.\psi) :: X$ is good below $|\sigma'|$ by the induction
hypothesis, so the claim for fuel $n$ gives
$\sigma' \in \fsem{\psi[\mu Y.\psi/Y]} = \fsem{\mu Y.\psi}$.

\begin{theorem}[\Th{Lemma~4.30}]\label{thm:vc}
  Assume that every procedure conforms with its contract, and let $\phi$
  be closed. Let $T'$ be $T$ with normalized bodies. If, with respect to
  $T'$, every $g$ that satisfies $\tc_1(\norm(S))$ also satisfies
  $\ftc^n_1(\norm(S), \phi, [\,])$, then $\Str\sem{S} \subseteq
  \fsem{\phi}$.
\end{theorem}

This holds for every statement $S$ and every procedure table $T$. Its
proof combines Lemmas~\ref{lem:eff-exists}, \ref{lem:tc} and
\ref{lem:norm} with Theorem~\ref{thm:ftc}, as in the thesis. The
formalization also proves the thesis's version, for $S$ and the procedure
bodies already in normal form, with $\tc_1(S) \Rightarrow
\ftc^n_1(S, \phi, [\,])$.

\section{The thesis's distribution function}\label{sec:distr-thesis}

This section shows that \Th{Theorem~4.29} and \Th{Lemma~4.30} fail when
$\ftc$ uses $\distr$ as in \Th{Definition~4.23}, which also distributes
chop over conjunction.

The thesis writes chop sequences without brackets: in \Th{Example~4.25},
$\distr((\Id \lor \Id\chop\Sb{x-1}{x})\chop\Sb{x+1}{x})$ is
$(\Id\chop\Sb{x+1}{x}) \lor (\Id\chop\Sb{x-1}{x}\chop\Sb{x+1}{x})$, read as
a flat sequence. HOL4 formulas are trees, so the formalization makes the
regrouping explicit.

\begin{definition}[\Th{Definitions~4.23 and~4.26}]
  \label{def:distr-thesis}
  $\distr_{\mathrm{FS}}(\phi)$ is defined as
  in \Th{Definition~4.23}, with the case for conjunction, plus the case
  \[
    \distr_{\mathrm{FS}}((\phi_1\chop\phi_2)\chop\phi_3) =
    \distr_{\mathrm{FS}}(\phi_1\chop(\phi_2\chop\phi_3)).
  \]
  A formula $\phi$ is \emph{normalised} if $\phi =
  \distr_{\mathrm{FS}}(\mathrm{expand}(\phi_0))$ for some $\phi_0$
  (Definition~\ref{def:expand}).
  $\ftc_{\mathrm{FS}}$ is $\ftc$ (Definition~\ref{def:ftc}) with
  $\distr_{\mathrm{FS}}$ in place of $\distr$.
\end{definition}

The added case preserves the semantics, since chop is associative, and it
only applies to a chop whose left operand is a chop. Let
$E = \mathrm{expand}(\mathit{true}) = \mathit{true} \lor
\mathit{true}\chop\mu X'.(\mathit{true} \lor \mathit{true}\chop X')$, so
$E$ describes the remaining statements of any length.

\begin{theorem}\label{thm:ftc-thesis}
  Let $T$ be empty, let $\Pre$ and $\Post$ be $\mathit{true}$, let
  $S = \skipS; (\skipS; \skipS)$, and let
  $\phi = \distr_{\mathrm{FS}}(\mathrm{expand}((\mu Y.\,\Id \land
  \Id\chop\Id)\chop\mathit{true})) = (\mu Y.\,\Id \land \Id\chop\Id)\chop E$.
  Then $\phi$ is closed and normalised, and $\fsem{\phi} = \emptyset$, but
  every $g$ that satisfies $\tc_1(S)$ satisfies
  $\ftc_{\mathrm{FS},1}^{10}(S, \phi, [\,])$. Hence
  \begin{itemize}
  \item \Th{Theorem~4.29} fails for $\ftc_{\mathrm{FS}}$: for a state $s$,
    $s\,s\,s\,s \eff_S s\,s\,s\,s$, and the constant sequence $g = s$
    satisfies $\ftc_{\mathrm{FS},1}^{10}(S, \phi, [\,])$, but
    $s\,s\,s\,s \notin \fsem{\phi}$;
  \item \Th{Lemma~4.30} fails for $\ftc_{\mathrm{FS}}$:
    $\tc_1(S) \Rightarrow \ftc_{\mathrm{FS},1}^{10}(S, \phi, [\,])$ is
    valid, but $s\,s\,s\,s \in \Str\sem{S} \setminus \fsem{\phi}$.
  \end{itemize}
\end{theorem}

The formula $\Id \land \Id\chop\Id$ has no traces, since $\Id$ has traces
of length~2 only and $\Id\chop\Id$ of length~3 only. $\ftc$ unfolds the
fixed point and distributes:
\[
  \distr_{\mathrm{FS}}((\Id \land \Id\chop\Id)\chop E) =
  (\Id\chop E) \land (\Id\chop(\Id\chop E)).
\]
It then checks the conjuncts separately. The first matches $\Id$ with the
first $\skipS$ and $E$ with the other two; the second matches
$\Id\chop\Id$ with the first two. The trace $s\,s\,s\,s$ is in both
conjuncts, but they split it at different positions, so it is not in
$(\Id \land \Id\chop\Id)\chop E$.

Without the regrouping, $\distr$ gives the second conjunct as
$(\Id\chop\Id)\chop E$, and $\ftc$ is false on it, because $\pred$ is false
on a chop. We expect that $\ftc$ with $\distr$ read this way is sound: then
$\ftc$ splits both conjuncts at the same statement boundaries. This is not
formalized. Without the regrouping, however, $\ftc$ cannot handle
normalised formulas such as those of \Th{Examples~4.31 and~4.32}.

\section{Examples}\label{sec:examples}

The examples check that the corrected definitions apply as the thesis
intends, and that the results are not vacuous. Chop sequences are
bracketed to the right. In \Th{Examples~3.1, 4.31 and~4.32}, the
procedure table is
\[
  T = m\{\ite{x \le 0}{x := 0}{(x := x - 1; \call{m})}\},
\]
and the precondition $\Pre_m$ is $\mathit{true}$.

\begin{lemma}\label{lem:m-conf}
  Every $\sigma \in \Str\sem{\call{m}}$ ends in a state where $x = 0$ and
  every other variable has its value in the first state. Hence $m$
  conforms with every contract whose postcondition follows from this.
\end{lemma}

The proof shows that the set of such traces is a pre-fixed point of the
transformer whose least fixed point interprets the procedures. The thesis
justifies conformance informally; its procedure contract condition is not
formalized.

\begin{theorem}[\Th{Example~3.1}]\label{thm:ex31}
  Let $S = y := x; \call{m}; y := y - x$ and $\Post_m = (x = 0 \land y =
  y_0)$, and let
  \[
    \phi = \Sb{x}{y}\chop\Id\chop(\mu X_m.\,\psi)\chop\Id, \qquad
    \psi = \Id\chop(\Sb{0}{x} \lor \Sb{x-1}{x}\chop\Id\chop X_m).
  \]
  Then $\tc_1(S) \Rightarrow \ftc^{20}_1(S, \phi, [\,])$ is valid, and
  hence $\Str\sem{S} \subseteq \fsem{\phi}$ by Theorem~\ref{thm:vc}.
  Moreover, $s^6 \in \Str\sem{S}$ for the state $s$ that maps every
  variable to $0$, and some trace of $y := x; \call{m}; y := y + 1$ is not
  in $\fsem{\phi}$.
\end{theorem}

$\ftc$ follows the shape of $S$. The call matches $\Id\chop\mu X_m.\psi$,
and the side condition checks $S_m$ against $\psi[\mu X_m.\psi/X_m]$, in
which the recursive call is discharged because $(m, \mu X_m.\psi) \in X$. The
trace of the variant is not in $\fsem{\phi}$ because every trace in
$\fsem{\phi}$ ends with two equal states.

\Th{Examples~4.31 and~4.32} use the specification
\[
  \phi' = \mathit{true}\chop(\mu X_m.\,\mathit{true}\chop(\Sb{0}{x} \lor
  X_m))\chop\Id,
\]
whose state formulas stand for any number of steps. $\ftc$ matches a state
formula with exactly one step, so $\phi'$ must be normalised as
$\distr(\mathrm{expand}(\phi'))$. After unfolding the fixed point, the
recursion variable for the call sits in the disjunction $\Sb{0}{x} \lor
X_m$ on the right of a chop, and the case for a call needs
$\phi_1\chop\mu X_m.\psi$. The thesis's computation of $\ftc(S, \phi)$ in
\Th{Example~4.31} is approximate and skips this step; the corrected
$\distr$ (Definition~\ref{def:distr}) makes it.

\begin{theorem}[\Th{Examples~4.31 and~4.32}]\label{thm:ex431}
  Let $\phi = \distr(\mathrm{expand}(\phi'))$.
  \begin{itemize}
  \item Let $S$ be as in Theorem~\ref{thm:ex31} and $\Post_m = (x = 0)$.
    Then $\tc_1(S) \Rightarrow \ftc^{30}_1(S, \phi, [\,])$ is valid, and
    hence $\Str\sem{S} \subseteq \fsem{\phi'}$.
  \item Let $S' = x := -x; z := 0; y := x; \call{m}; y := y - (x + z)$.
    With $\Post_m = (x = 0)$, some $g$ satisfies $\tc_1(S')$ but not
    $\ftc^n_1(S', \phi, [\,])$, for every $n$. With $\Post_m = (x = 0
    \land z = z_0)$, $\tc_1(S') \Rightarrow \ftc^{30}_1(S', \phi, [\,])$ is
    valid, and hence $\Str\sem{S'} \subseteq \fsem{\phi'}$.
  \end{itemize}
\end{theorem}

The conclusions use $\fsem{\phi} = \fsem{\phi'}$
(\Th{Proposition~4.27}). The failure for the weaker contract holds for
every fuel, so it is not an artifact of fuel: after the call, the
contract does not constrain $z$, which the last assignment needs. $S'$
satisfies $\phi'$ nonetheless, as the stronger contract shows.

\begin{proposition}[\Th{Examples~4.12, 4.15, 4.18 and~4.25}]
  \label{prop:ex-aux}\leavevmode
  \begin{itemize}
  \item \Th{Example~4.12}: for $S = (\ite{x = 0}{x := 1}{((\ite{x <
    0}{\skipS}{x := -x}); \skipS)}); x := x + 1$, $\norm(S)$ is the
    statement $S'$ of the thesis, and it is in normal form.
  \item \Th{Example~4.15}: for $S = (\ite{y = 0}{x := 1}{\skipS}); x :=
    2$, every $g$ that satisfies $\tc_1(S)$ has $g(1)(y) \neq 0$ and
    $g(1)(x) = 2$, while some $g$ with $g(1)(y) = 0$ satisfies
    $\tc_1(\norm(S))$.
  \item \Th{Example~4.18}: for $\phi = \mu X.\,\Id\chop X \lor ((x = 10)
    \land \Sb{x+1}{x})$, $\pred^1_i(\phi)$ holds iff $g(i)(x) = 10$ and
    $g(i+1) = g(i)[x \mapsto g(i)(x) + 1]$, and the trace
    $[x \mapsto 10]\,[x \mapsto 11]$ is in $\fsem{\phi}$.
  \item \Th{Example~4.25}: $\distr((\Id \lor \Id\chop\Sb{x-1}{x})\chop
    \Sb{x+1}{x}) = \Id\chop\Sb{x+1}{x} \lor
    \Id\chop(\Sb{x-1}{x}\chop\Sb{x+1}{x})$.
  \end{itemize}
\end{proposition}

The thesis says that $\tc(S)$ of \Th{Example~4.15} is unsatisfiable
regardless of $y$; it is satisfiable when $y \neq 0$ and $x = 2$
initially.

\section{Incompleteness}\label{sec:incomplete}

The converse of Theorem~\ref{thm:vc} fails: a statement may satisfy its
specification while $\tc(S) \Rightarrow \ftc(S, \phi)$ is not valid
(\Th{Section~6.3}). The formalization shows this for \Th{Examples~6.5
and~6.6}. In both, the failure holds for every contract that the
procedures conform with, so a stronger contract does not help, and for
every fuel.

\begin{theorem}[\Th{Example~6.5}]\label{thm:ex65}
  Let
  \[
    T = m\{\ite{x > 0}{(y := 2 \cdot y; x := x - 1; \call{m})}{x :=
    0}\},
  \]
  $S = x := 5; y := 1; \call{m}$ and $\phi = \mathit{true}\chop\Sb{8}{y}
  \chop\mathit{true}$. Then $\Str\sem{S} \subseteq \fsem{\phi}$. For every
  contract that $m$ conforms with, some $g$ satisfies $\tc_1(S)$ but not
  $\ftc^n_1(S, \distr(\mathrm{expand}(\phi)), [\,])$, for every $n$.
\end{theorem}

The thesis writes $\Sb{2^3}{y}$; arithmetic expressions have no
exponentiation, so the formalization writes $8$. Every trace of $S$ sets
$y$ to $8$ in the third iteration of $m$, when $y = 4$. The proof shows
that the traces of $\call{m}$ that start with $(x, y)$ equal to $(5, 1)$,
$(4, 2)$ or $(3, 4)$ have such a step, by the pre-fixed-point argument of
Lemma~\ref{lem:m-conf}. The verification condition fails because $\ftc$
can check $\Sb{8}{y}$ only with $\pred$, at one of the steps of $x := 5$,
$y := 1$ and the duplicated first state of the call, and none of them
sets $y$ to $8$. The only fixed point that can specify the call is the
one that $\mathrm{expand}$ introduces for the last $\mathit{true}$, which
does not mention $y$: the contract describes only the first and last
states of the call. The witness $g$ needs a final state for the call that
satisfies the contract; it is the last state of a trace of $\call{m}$,
which terminates from every state with $x \ge 0$.

\begin{theorem}[\Th{Example~6.6}]\label{thm:ex66}
  Let $T = m\{x := x + 1; \call{n}\}, n\{y := y + 1\}$, $S = \call{m}$ and
  $\phi = \Id\chop\Sb{x+1}{x}\chop\Id\chop\Sb{y+1}{y}$. Then
  \begin{itemize}
  \item $\Str\sem{S} = \fsem{\phi}$ and $\distr(\mathrm{expand}(\phi)) =
    \phi$;
  \item $\ftc^n_1(S, \phi, [\,])$ is false for every $g$, $n$ and
    contract, while $\tc_1(S)$ is satisfiable for every contract that $m$
    and $n$ conform with;
  \item for $\phi_\mu = \Id\chop\mu X_m.\,(\Sb{x+1}{x}\chop\Id\chop\mu
    X_n.\,\Sb{y+1}{y})$, $\fsem{\phi_\mu} = \fsem{\phi}$, and with
    $\Pre_m = \Pre_n = \mathit{true}$ and any postconditions, $\tc_1(S)
    \Rightarrow \ftc^{10}_1(S, \phi_\mu, [\,])$ is valid.
  \end{itemize}
\end{theorem}

So $\Str\sem{S} \subseteq \fsem{\phi}$ holds, with equality: $\phi$ is the
strongest trace formula of $S$ with the fixed points unfolded. $\ftc$
fails because the case for a call needs $\phi_1\chop\mu Y.\psi$, and
$\phi$ has no fixed point. With the fixed points of $\phi_\mu$, as the
thesis suggests, the verification condition holds, and
Theorem~\ref{thm:vc} gives $\Str\sem{S} \subseteq \fsem{\phi_\mu}$.

The formalized $\ftc$ accepts any fixed point for a call
(Definition~\ref{def:ftc}), while the thesis's \Th{Definition~4.28}
suggests that the recursion variable must be $X_m$. The theorems above are
about the more permissive $\ftc$. For the thesis's version, the
verification conditions of both examples fail for the same reasons; in
\Th{Example~6.5}, the fixed point from $\mathrm{expand}$ then cannot
specify the call at all. This is not formalized.

\section{What is not formalized}

This formalization stops at \Th{Lemma~4.30}, with the encodings as
predicates on state sequences. The first-order logic of
\Th{Section~4.1}, the procedure contract condition and
\Th{Lemmas~4.35 and~4.36}, a conversion to first-order verification
conditions with end-to-end soundness, and the evaluation examples of
\Th{Section~5.4} are formalized on top of it, as described in a separate
note (\texttt{doc/fol.tex}). Not formalized are the other examples of
the thesis, among them those of \Th{Sections~6.1 and~6.2}, which
illustrate alternative designs of the calculus.

\begin{thebibliography}{GH}
\bibitem[GH]{GH} D.~Gurov and R.~H\"ahnle. \emph{An Expressive Trace Logic
  for Recursive Programs}.
  \url{https://doi.org/10.48550/arXiv.2411.13125}.
\bibitem[FS]{FS} F.~Stevanovic. \emph{Automated Verification of Finite-trace
  Properties for Recursive Programs}. Master's thesis, KTH Royal Institute
  of Technology, 2025.
  \url{https://urn.kb.se/resolve?urn=urn\%3Anbn\%3Ase\%3Akth\%3Adiva-368044}.
\end{thebibliography}

\end{document}
